\begin{summarybox}

	\textbf{\textcolor{CSummary}{一句话核心}}

	\textbf{
		若事件 $A$ 与 $B$ 相互独立，则
		\[
			P(AB)=P(A)P(B),
		\]
		其中 $AB=A\cap B$ 表示两个事件同时发生。
	}

	\medskip

	\textbf{\textcolor{CSummary}{使用条件}}

	\begin{itemize}[leftmargin=*, itemsep=0.5em, topsep=0.4em]
		\item \textbf{条件 1（事件独立）：}
		      题目明确说明事件 $A$ 与 $B$ 相互独立，或已通过条件验证二者独立。

		\item \textbf{条件 2（多个事件）：}
		      使用多个事件的乘积公式时，要求这些事件相互独立，而不能只凭“两两独立”直接判断。
	\end{itemize}

	\medskip

	\textbf{\textcolor{CSummary}{关键提醒}}

	\begin{itemize}[leftmargin=*, itemsep=0.5em, topsep=0.4em]
		\item \textbf{易错点（忽视独立）：}
		      未确认独立性，就直接套用
		      \[
			      P(AB)=P(A)P(B).
		      \]

		\item \textbf{检查项（乘积检验）：}
		      若能够求出 $P(AB)$、$P(A)$ 与 $P(B)$，则可通过检验
		      \[
			      P(AB)=P(A)P(B)
		      \]
		      是否成立，判断 $A,B$ 是否相互独立。

		\item \textbf{补充项（条件概率表述）：}
		      当 $P(A)>0$ 时，事件 $A$ 与 $B$ 相互独立等价于
		      \[
			      P(B\mid A)=P(B).
		      \]
		      当 $P(A)=0$ 时，$P(B\mid A)$ 没有定义，不能使用这一条件概率表述。
	\end{itemize}

\end{summarybox}